\subsubsection{Constant Global Time Stepping}
\label{SEC:Time:Algorithms:Constant}

Constant global time stepping is exactly as it sounds: a constant time step
size, taken from the input \inp[nmlTimestep::Timestep]{Timestep}, is globally
used by all cells to update their solutions in time.
%
This method is used by setting the input
\inp[nmlTimestep::TimestepType]{Timestep\_Type}=0.
%
In general, the time step size is constant in that the value of
\inptext{Timestep} is used for every solver iteration with a
\inp[nmlTimestep::CFL]{CFL} of 1.0.
%
However, if CFL ramping is being used, the time step size will change based on
the ramping schedule.
%
Once the ramping schedule finishes, \inptext{CFL} will revert to a value of
$1.0$.
%
%Once the ramping schedule finishes, the time step size will go back to being
%constant with a value of \[\Delta t = \inp[nmlTimestep::Timestep]{Timestep}
%\times \inp[nmlTimestep::CFLRampingEnd]{CFL\_Ramping\_End} \]

%%%%%%%%%%%%%%%%%%%%%%%%%%%%%%%%%%%%%%%%%%%%%%%%%%%%%%%%%%%%%%%%%%%%%%%%%%%%%%%%

Being a global time stepping method, this option is suitable for time-accurate
simulations like LES.
%
The difficulty in using a constant global time step is that the maximum
allowable time step for a simulation is problem dependent, so it can only really
be used when you have an idea of what time step size is needed for a problem.
%
The primary use for constant time stepping would be for running a followup
simulation using the restart file of a simulation that used a CFL-based global
time step in which all initial transients have been removed from the flow.
%
However, continuing the use of CFL-based global time stepping for this followup
situation is just as good as, if not better than, switching to a constant time
step size, since the CFL-based approach can reduce the time step size if
something develops within the simulation that requires a smaller time step.

%%%%%%%%%%%%%%%%%%%%%%%%%%%%%%%%%%%%%%%%%%%%%%%%%%%%%%%%%%%%%%%%%%%%%%%%%%%%%%%%

